\section{A physical Gram matrix and the purity obstruction}
\label{sec:ext-gram}

The general spectral formula of Section~\ref{sec:formula} treats every value of the
unilateral optimum, including sub-classical values and singular pencils. Physical
positivity allows a much simpler formula throughout the entire region in which
unilateral CHSH violation is possible, and we present it first. The simplification also yields
polynomial activation tests and useful probability tradeoffs.

Continue to use the canonical data in~\eqref{eq:data}, and write
\begin{equation}
 \alpha=\tr A-s,\qquad m=\frac{\alpha-1}{2},\qquad H_0=K+iC.
 \label{eq:ext-alpha}
\end{equation}
Here and below $\alpha$ denotes this purity invariant. The complex matrix
$H_0$ is Hermitian. We use $z^\dagger z$ for the Hermitian squared norm and
$z^Tz$ for the complex bilinear quadratic form; the distinction will matter
at the CHSH threshold.

\begin{lemma}[Spin-flip Gram identity]
\label{lem:ext-gram}
Let
\[
 \widetilde\rho_c=(\sigma_y\otimes\sigma_y)\rho_c^T
                    (\sigma_y\otimes\sigma_y),\qquad
 S_i=\sigma_i\otimes I,
 \qquad
 G_{ij}=\tr(\rho_c S_i\widetilde\rho_c S_j).
\]
Then
\begin{equation}
 \boxed{\,2G=K-mI+iC\succeq0.\,}
 \label{eq:ext-gram}
\end{equation}
In particular $K\succeq mI$, or equivalently
\begin{equation}
 \tr A+s\le 1+2\lambda_{\min}(A).
 \label{eq:ext-canonical-physicality}
\end{equation}
For real vectors $n,v$ and $z=n+iv$,
\begin{equation}
 z^\dagger(K-mI+iC)z
 =n^T(K-mI)n+v^T(K-mI)v+2v^TCn.
 \label{eq:ext-real-gram}
\end{equation}
\end{lemma}

\begin{proof}
The spin flip is positive semidefinite, since it is a unitary conjugate of
the transpose of a positive semidefinite matrix. For $Z=\sum_jz_jS_j$,
\[
 z^\dagger Gz
 =\tr(\rho_c Z^\dagger\widetilde\rho_c Z)
 =\norm{\widetilde\rho_c^{1/2}Z\rho_c^{1/2}}_F^2\ge0.
\]
For completeness the Pauli coefficient calculation is as follows. The
identity--identity contribution to $G_{ij}$ is $\delta_{ij}/4$; the
Bob-marginal contribution is $-s\delta_{ij}/4$; and the
correlation--correlation contribution is
$(2A_{ij}-\tr A\,\delta_{ij})/4$. The two remaining cross terms are
\[
 \frac{i}{4}\sum_k\epsilon_{ikj}d_k
 -\frac{i}{4}\sum_k\epsilon_{kij}d_k
 =-\frac{i}{2}\sum_k\epsilon_{ijk}d_k
 =\frac{i}{2}C_{ij}.
\]
Their sum is
$G=\tfrac12 A+\tfrac14(1-s-\tr A)I+\tfrac i2 C
 =\tfrac12(K-mI+iC)$.
Expanding $z=n+iv$ and using $C^T=-C$ gives
\eqref{eq:ext-real-gram}. Taking $v=0$ gives the real-matrix inequality
$K\succeq mI$.
\end{proof}

The bare implication from CHSH violation to a negative conditional
R\'enyi-2 entropy is established in the literature~\cite{horodecki1996}.
The next proof records its precise quantitative form and its covariance
under unilateral filtering, because these are the ingredients needed for
the spectral simplification. The plane-purity identity
\eqref{eq:ext-plane-purity} and the determinant covariance
\eqref{eq:ext-delta-covariance}, with the roles of the two parties
exchanged, also appear in the proof of Conjecture~1 of Milne, Jennings,
Jevtic and Rudolph recorded in the Omega Institute's public
repository~\cite{dossier} (Section~\ref{sec:prior}); they are not claimed
as new here.

\begin{proposition}[Purity obstruction and determinant covariance]
\label{prop:ext-purity}
For a possibly unnormalized Hermitian two-qubit matrix $X$, define
\[
 \Delta_B(X)=\tr X^2-\tr[(\tr_A X)^2].
\]
For every complex $2\times2$ matrix $F$,
\begin{equation}
 \Delta_B((F\otimes I)X(F^\dagger\otimes I))
 =|\det F|^2\Delta_B(X).
 \label{eq:ext-delta-covariance}
\end{equation}
For every normalized two-qubit state,
\begin{equation}
 B(\rho)^2\le1+2\Delta_B(\rho).
 \label{eq:ext-purity-bound}
\end{equation}
Consequently Alice-only CHSH violation requires $\Delta_B(\rho)>0$.
For a canonical state and a filter with effect proportional to
$I+u\cdot\sigma$,
\begin{equation}
 B((\rho_c)_u)^2\le1+m(1-\norm u^2).
 \label{eq:ext-filter-purity}
\end{equation}
Thus $B_A^*>1$ implies $\alpha>1$ and $K\succ0$.
\end{proposition}

\begin{proof}
Expand $X=\tfrac12\sum_{\nu=0}^3X_\nu\otimes\sigma_\nu$, with Hermitian
$2\times2$ coefficients $X_\nu$. Pauli orthogonality and
$\tr Y^2-(\tr Y)^2=-2\det Y$ give
\[
 \Delta_B(X)
 =\frac12\sum_\nu[\tr(X_\nu^2)-(\tr X_\nu)^2]
 =-\sum_\nu\det X_\nu.
\]
Every coefficient transforms as $X_\nu\mapsto FX_\nu F^\dagger$,
which proves~\eqref{eq:ext-delta-covariance}, including singular $F$.

For the quantitative bound, let $(a,b,T)$ be the Bloch data of $\rho$,
let $n$ be a real unit vector, and conjugate the full spin flip
$\widetilde\rho$ by $(n\cdot\sigma)\otimes I$ to obtain the positive
matrix $\widetilde\rho_n$. Conjugation on Alice acts on Bloch vectors
by $2nn^T-I$, so Pauli orthogonality gives
\[
 4\tr(\rho\widetilde\rho_n)
 =1+\norm a^2-\norm b^2-\tr(TT^T)
       -2(a\cdot n)^2+2n^TTT^Tn.
\]
Together with
$4\Delta_B(\rho)=\norm a^2-\norm b^2+\tr(TT^T)-1$, this proves
the exact identity
\begin{equation}
 \tr(TT^T)-n^TTT^Tn+(a\cdot n)^2
       +2\tr(\rho\widetilde\rho_n)
 =1+2\Delta_B(\rho).
 \label{eq:ext-plane-purity}
\end{equation}
Both of the last two terms on the left are nonnegative. Minimizing
$n^TTT^Tn$ proves~\eqref{eq:ext-purity-bound}. Equality requires a
maximizing CHSH plane whose omitted direction has $a\cdot n=0$ and
$\tr(\rho\widetilde\rho_n)=0$.

If a filter succeeds with probability $p$, covariance and normalization
give
\[
 B(\rho_F)^2
 \le1+\frac{2|\det F|^2}{p^2}\Delta_B(\rho).
\]
In canonical coordinates,
$\Delta_B(\rho_c)=(\alpha-1)/4$, and the effect $I+u\cdot\sigma$
has normalization one and determinant $1-\norm u^2$.
This proves~\eqref{eq:ext-filter-purity}. The same inequality with
$\alpha\le1$ excludes every violation, including any limiting violation.
If $\alpha>1$, Lemma~\ref{lem:ext-gram} gives
$K\succeq mI\succ0$.
\end{proof}

For the original Bloch data, put $h=1-\norm a^2>0$. Whitening has
$|\det W|^2=1/h$, so covariance gives the useful identities
\begin{equation}
 \boxed{\quad
 \alpha-1=\frac{\norm{T}_F^2+\norm a^2-\norm b^2-1}{h},
 \qquad
 \alpha=\frac{\norm{T}_F^2-\norm b^2}{h}.
 \quad}
 \label{eq:ext-raw-alpha}
\end{equation}
Here $\norm{T}_F$ is the Frobenius norm. Thus the purity guard is
already a quadratic polynomial inequality in the original state
coordinates.

\section{The smallest eigenvalue is automatically feasible}
\label{sec:ext-shortcut}

\begin{theorem}[A single smallest eigenvalue in the violation-relevant region]
\label{thm:ext-shortcut}
For a physical canonical state with $\alpha>1$, define
$M=A+CK^{-1}C$. Then
\begin{equation}
 \boxed{\qquad (B_A^*)^2=\tr A-\lambda_{\min}(M).\qquad}
 \label{eq:ext-shortcut}
\end{equation}
Every unit vector $n$ in the minimum eigenspace gives an optimal
invertible filter:
\begin{equation}
 v=-K^{-1}Cn,\qquad u=-n\times v,\qquad
 F\ \text{proportional to}\ \sqrt{I+u\cdot\sigma}\,W.
 \label{eq:ext-shortcut-filter}
\end{equation}
In particular, $n\cdot v=0$ and $\norm v<1$ hold automatically.
All optimal pairs in the compact reduction are precisely the pairs in
\eqref{eq:ext-shortcut-filter}.

The same conclusion holds under the algebraic hypothesis
$K+iC\succ0$, even when $\alpha\le1$.
\end{theorem}

\begin{proof}
The Gram identity gives $H_0\succeq mI\succ0$ when $\alpha>1$.
Under either hypothesis in the theorem, the real representation
\[
 \begin{pmatrix}K&-C\\ C&K\end{pmatrix}
\]
is positive definite. Its Schur complement is
$L=K+CK^{-1}C=M-sI\succ0$. Let
$\delta=\lambda_{\min}(L)>0$, choose $\norm n=1$ with $Ln=\delta n$,
and put $v=-K^{-1}Cn$. Then
\begin{equation}
 Kn-Cv=\delta n,\qquad Kv+Cn=0.
 \label{eq:ext-paired-equations}
\end{equation}
Taking the inner product of the first equation with $v$ and of the
second with $n$ gives $\delta n\cdot v=0$. Hence $n\cdot v=0$.

The crucial norm estimate is not a generic property of a stationary
eigenvector. It follows here from the following exact identity:
\begin{equation}
 \boxed{\,v^TLv=\delta-\delta^2n^TK^{-1}n.\,}
 \label{eq:ext-rayleigh-identity}
\end{equation}
Indeed $Cv=(K-\delta I)n$, and taking the inner product of the
first equation in~\eqref{eq:ext-paired-equations} with $n$ yields
$v^TKv=n^TKn-\delta$. Therefore
\begin{align*}
 v^TLv
 &=v^TKv-(Cv)^TK^{-1}(Cv)\\
 &=n^TKn-\delta
   -n^T(K-\delta I)K^{-1}(K-\delta I)n\\
 &=\delta-\delta^2n^TK^{-1}n.
\end{align*}
The minimum Rayleigh bound for $L$ now gives
\begin{equation}
 \norm v^2\le1-\delta n^TK^{-1}n<1.
 \label{eq:ext-feasible-margin}
\end{equation}
Thus the constructed pair belongs to the interior of the exact compact
domain~\eqref{eq:domain}.

For the matching upper bound, complete the square in~\eqref{eq:f}:
\begin{equation}
 f(n,v)=\tr A-n^TMn
 -(v+K^{-1}Cn)^TK(v+K^{-1}Cn).
 \label{eq:ext-square-completion}
\end{equation}
The last term is nonnegative, and $n^TMn\ge\lambda_{\min}(M)$.
The constructed pair gives equality in both inequalities. Because
$K\succ0$, equality also characterizes every optimizing pair: $n$
must lie in the minimum eigenspace and $v=-K^{-1}Cn$.
\end{proof}

\begin{corollary}[Strict boundary gap and a uniform probability guarantee]
\label{cor:ext-boundary-gap}
Whenever $H_0\succ0$, the optimum strictly exceeds the rank-one
boundary value: $(B_A^*)^2>E$, even if the optimum is at most one.
Writing $\delta=\lambda_{\min}(M)-s>0$, every canonical optimal
effect obeys
\begin{equation}
 \norm u^2\le1-\frac{\delta}{\lambda_{\max}(K)},\qquad
 p_{\max}\ge
 \frac1{1+\sqrt{1-\delta/\lambda_{\max}(K)}}>\frac12.
 \label{eq:ext-probability-margin}
\end{equation}
If the minimum eigenvalue of $M$ is simple, the positive optimal
filter is unique up to scale. An arbitrary output unitary remains
irrelevant to the CHSH value.
\end{corollary}

\begin{proof}
Every equality pair in~\eqref{eq:ext-square-completion} satisfies
the strict estimate~\eqref{eq:ext-feasible-margin}. No boundary pair
can therefore attain the global maximum. Compactness of the boundary
gives the strict gap. Moreover
$n^TK^{-1}n\ge1/\lambda_{\max}(K)$ gives the first inequality in
\eqref{eq:ext-probability-margin}; its radicand is nonnegative because
$M-sI=K+CK^{-1}C\preceq K$.
The probability statement follows by maximal Kraus rescaling.
For a simple eigenvalue, $n$ is unique up to sign, and simultaneous
sign change of $n,v$ leaves $u=-n\times v$ fixed.
The equality discussion following~\eqref{eq:attainment} shows that
this reconstruction recovers the original positive-filter part of
every finite optimum.
\end{proof}

\begin{corollary}[Sharp purity ceiling and its equality cases]
\label{cor:ext-ceiling}
If $\alpha>1$, then
\begin{equation}
 (B_A^*)^2\le\frac{1+\alpha}{2}.
 \label{eq:ext-ceiling}
\end{equation}
Equality holds exactly when
\begin{equation}
 \ker(K-mI)\cap\ker C\ne\{0\}.
 \label{eq:ext-ceiling-equality}
\end{equation}
In that case the canonical state itself is optimal. Every state of rank
at least three satisfies~\eqref{eq:ext-ceiling} strictly.
\end{corollary}

\begin{proof}
Apply~\eqref{eq:ext-filter-purity} to the optimal filter supplied by
Theorem~\ref{thm:ext-shortcut}. Since $m>0$, equality forces its
canonical effect to have $u=0$. The omitted direction $n$ then satisfies
$n^T(K-mI)n=0$. Positivity of $K-mI+iC$ implies that this matrix
annihilates $n$; its real and imaginary parts yield
$(K-mI)n=0$ and $Cn=0$. Conversely, such a unit vector gives
$f(n,0)=\alpha-m=(1+\alpha)/2$.

For the rank statement, a real nonzero $n$ gives an invertible
$Z=(n\cdot\sigma)\otimes I$. If $r=\rank\rho_c\ge3$, then
\[
 \rank(\widetilde\rho_c^{1/2}Z\rho_c^{1/2})\ge2r-4>0.
\]
Thus its squared Frobenius norm cannot vanish, excluding a real
nonzero vector in $\ker G$. More generally, the same argument excludes
any $z\in\ker G$ with $z^Tz\ne0$, since
$\det(z\cdot\sigma)=-z^Tz$.
\end{proof}

\subsection{Why positivity of \(K\) alone is insufficient}

The purity or positive-Gram hypothesis in
Theorem~\ref{thm:ext-shortcut} cannot be replaced by $K\succ0$.
Consider the full-rank canonical state
\begin{equation}
 \rho_{\mathrm{guard}}=\frac1{40}
 \begin{pmatrix}
 19&0&0&7\\0&1&0&0\\0&0&7&0\\7&0&0&13
 \end{pmatrix}.
 \label{eq:ext-counterexample}
\end{equation}
Its parameters are
\[
 \begin{aligned}
 c&=(0,0,3/10),&
 D&=\diag(7/20,-7/20,3/5),\\
 K&=\diag(13/400,13/400,27/100)\succ0,&
 \alpha&=103/200.
 \end{aligned}
\]
The unrestricted smallest-eigenvalue expression gives
\[
 \tr A-\lambda_{\min}(M)=\frac{7693}{5200}>1,
 \qquad \norm v=\frac{72}{13}>1.
\]
The spurious minimum eigenspace is the entire transverse plane, and
every unit vector in that plane has the same infeasible norm.
The only remaining interior eigenvalue gives $49/200$.
For the boundary radius, writing $z=u_3$ yields a convex quadratic
on $[-1,1]$ whose largest endpoint is $(3/10+3/5)^2=81/100$.
The general spectral theorem therefore gives the exact answer
\[
 \boxed{(B_A^*(\rho_{\mathrm{guard}}))^2=81/100.}
\]
This example explains why an unguarded cubic formula can incorrectly
predict a CHSH violation.

\subsection{A full-rank violating state with a circle of optimal filters}

The simple-eigenvalue hypothesis in the uniqueness statement cannot be
discarded, even for a full-rank state with a strictly violating optimum.
Take
\[
 c=(9/20,0,0),\qquad D=\diag(8/15,-2/3,3/4),
\]
or equivalently
\begin{equation}
 \rho_{\mathrm{circle}}=
 \begin{pmatrix}
 7/16&9/80&0&3/10\\
 9/80&1/16&-1/30&0\\
 0&-1/30&1/16&9/80\\
 3/10&0&9/80&7/16
 \end{pmatrix}.
 \label{eq:ext-circle-state}
\end{equation}
In the Bell-parity basis its two blocks are
$\left(\begin{smallmatrix}59/80&9/80\\9/80&7/240\end{smallmatrix}\right)$
and
$\left(\begin{smallmatrix}11/80&9/80\\9/80&23/240\end{smallmatrix}\right)$.
Their determinants are $17/1920$ and $1/1920$, respectively, and
their diagonals are positive, giving an exact certificate that
$\rho_{\mathrm{circle}}\succ0$.
Here
\[
 \begin{aligned}
 \alpha&=\frac{49}{45},\qquad d=(6/25,0,0),\\
 K&=\diag\left(\frac{59}{720},\frac{871}{3600},\frac9{25}\right),\\
 M-sI&=\diag\left(\frac{59}{720},\frac{59}{720},\frac{531}{4355}\right).
 \end{aligned}
\]
The minimum eigenspace is $\operatorname{span}(e_1,e_2)$.
For $n=(\cos\theta,\sin\theta,0)$,
$v=(0,0,-2\sin\theta/3)$, so the complete set of optimal canonical
effects is
\begin{equation}
 u=\frac23(\sin^2\theta,-\cos\theta\sin\theta,0),
 \qquad
 (u_1-\tfrac13)^2+u_2^2=\frac19,\quad u_3=0.
 \label{eq:ext-optimal-circle}
\end{equation}
Every point of this circle attains
$(B_A^*)^2=145/144>1$.
The circle includes the identity effect $u=0$, and all its points
have $\norm u\le2/3<1$. Thus equality cases can form a genuine
continuous family of distinct positive filters without loss of
rank or CHSH violation.
